\section{三、研究设计}\label{ux4e09ux7814ux7a76ux8bbeux8ba1}

\subsection{（一）样本选择与数据来源}\label{ux4e00ux6837ux672cux9009ux62e9ux4e0eux6570ux636eux6765ux6e90}

本文以2011至2024年沪深A股上市公司为样本。数据包括两个层面，年报全文用于提取数据要素利用指标，CSMAR数据库的财务、治理和市场数据用于构造因变量和控制变量。样本筛选依据中国证监会行业分类标准剔除J类金融保险业企业，剔除ST/*ST企业，剔除上市时间在观测年份内不足12个月的企业-年份观测。对连续变量采取1\%和99\%分位数的Winsorize缩尾处理。

变量的时间匹配遵循按报告期对齐的原则（Feng等，2024；Wei等，2025）：t年年报覆盖企业t年经营活动，PriceDelay基于t年日收益率计算，二者同属一个会计年度。考虑到英文审稿中可能提出的``用事后文本编码当年信息环境''的质疑，本文在第四部分除报告同期规格外，还将\(t-1\)期DU解释\(t\)期PriceDelay的滞后规格前置呈现，作为时间顺序更清晰的偏好识别规格。年报文本提取的时间范围为2010至2024年，其中2010年数据仅用于该滞后分析。

经过上述筛选，最终获得43,735个企业-年份观测，涵盖约4,950家独立企业。OLS基准回归和工具变量估计均基于该样本，扣除固定效应单例后进入回归的观测约43,700。DML估计因使用27个控制变量和交叉拟合，样本约35,700，在相应章节单独说明。

年报全文来源于上市公司向证监会提交的法定年报TXT格式文件。因变量PriceDelay基于CSMAR数据库的股票日收益率数据构造，控制变量和机制变量同样来自CSMAR。

\subsection{（二）数据要素利用的测度}\label{ux4e8cux6570ux636eux8981ux7d20ux5229ux7528ux7684ux6d4bux5ea6}

数据资产及其相关活动长期处于会计体系之外，传统财务指标无法直接观察（Coyle和Manley，2024），文本分析是目前获取相关信息的主要方式。本文沿数据价值链的不同环节构建了五维度关键词体系，共计73个关键词（见下表）。

\textbf{表1 数据要素利用关键词体系}

{\def\LTcaptype{} % do not increment counter
\begin{longtable}[]{@{}lll@{}}
\toprule\noalign{}
维度 & 关键词数 & 示例 \\
\midrule\noalign{}
\endhead
\bottomrule\noalign{}
\endlastfoot
数据存量 data\_stock & 15 & 大数据、数据库、数据中心 \\
数据开发能力 data\_dev & 19 & 数据挖掘、机器学习、数字化转型 \\
数据驱动应用 data\_app & 19 & 精准营销、智能推荐、风控模型 \\
数据价值变现 data\_value & 11 & 数据资产、数据交易、数据入表 \\
数据治理 data\_gov & 9 & 数据安全、数据隐私、数据合规 \\
\end{longtable}
}

这一分类方案参考既有数字化文本测度，并针对朱康（2025）采用的两维度体系进行了细化，新增数据开发能力、数据价值变现和数据治理三个维度，覆盖数据要素利用的主要环节。完整关键词清单见附录。

关键词频率的计算方法为：

\[DU_{kw} = \frac{\text{关键词总出现次数}}{\text{年报全文总字数}} \times 10000 \qquad (1)\]

即每万字中数据相关关键词的出现次数。文本处理采用多编码支持策略和长词优先匹配策略，以准确读取多种编码格式并避免短词误匹配。

关键词频率统计存在测量噪声：企业可能在年报中笼统地提及``大数据''\,``数字化''而并无实质利用，也可能在不同语境下使用相同关键词但利用深度迥异。为缓解这一问题，本文引入大语言模型语义评分作为第二套测度。

具体步骤如下：首先基于73个关键词定位年报中所有命中位置及其前后100字的上下文窗口；然后使用大语言模型对每个上下文进行0-3分的语义评分，判断该处关键词是否反映了实质性的数据要素利用（0=无实质内容，3=深度利用）；最后取企业-年份层面各上下文评分的众数作为该企业当年的LLM评分llm\_score，以抑制少数噪声上下文的拉动。评分使用Claude Haiku 4.5模型，温度参数设为0以确保输出确定性。附录进一步报告评分模板、0/1/2/3判分规则和真实上下文示例。

在此基础上构造三个衍生指标。第一，LLM二值化指标llm\_binary在llm\_score\textgreater=2时取1否则取0。第二，归档语义增强指标DU\_llm = ln(1 + kw\_total) × (llm\_score / 3)，将关键词总出现次数取对数后与归一化的语义质量相乘。第三，长度标准化语义深度指标DU\_llm\_lenstd = ln(1 + DU\_kw) × (llm\_score / 3)，以已做篇幅标准化的DU\_kw替换kw\_total，对语义深度进行更保守的口径校正。本文将DU\_kw明确作为基准处理变量和主要识别对象，将DU\_llm作为归档语义测度并行报告，并以DU\_llm\_lenstd检验语义深度在控制篇幅后的稳健性。

为识别更宽泛数字叙事与数据利用披露之间的边界，本文进一步构造GenericNarr，即八个最泛化数字叙事词------``数字化转型、数字化、智能化、信息化、人工智能、大数据、智能制造、智慧城市''------的每万字出现频次；并据此定义去泛化词指标DU\_kw\_strict，即从kw\_total中扣除上述八个词后的每万字频次。扩展分析还构造WashGap = z(DU\_kw) - z(DU\_llm\_lenstd)，刻画``提得多但语义深度不足''的广度---深度错配；同时定义BroadShallow虚拟变量，当企业当年DU\_kw处于年度上三分位且DU\_llm\_lenstd处于年度下三分位时取1，否则取0。

稳健性检验还使用两个备选指标：DU\_kw\_ln = ln(1 + DU\_kw)，即对频率进行对数变换以处理右偏分布；DU\_sub\_ln = ln(1 + DU\_sub)，其中DU\_sub为基于规则判断的实质利用次数。

在上述双测度之外，本文进一步构造四类增强口径，以检验测度能否从``泛化数字叙事''推进到``更接近数据资源实际使用''的表述。第一，DUevent\_raw和DUevent分别定义为规则识别的实质利用次数及其每万字标准化值。这里的规则识别基于关键词附近上下文中的动作词、具体业务对象和数值信息，substantive\_count越高，表示文本中越多表述具有``被用于具体经营或治理场景''的特征。第二，DUchain\_count定义为五个维度中被命中的维度数，取值为0至5；DUclosedloop为闭环式表述虚拟变量，当企业同时出现数据资源/存量类、应用/开发类以及治理/价值实现类表述时取1，否则取0。第三，DUcore\_raw和DUcore聚焦数据驱动应用、价值实现和治理合规三类更接近经营与制度闭环的词项强度，用于形成更窄的支持性口径。第四，DUkw\_mda为MD\&A章节内关键词频率，衡量数据利用表述是否更集中于经营讨论部分。本文将这些增强指标明确定位为支持性测度：其目的在于检验``更完整、更嵌入''的数据利用信息是否同样具有资本市场定价含量，而非替换DU\_kw这一基准处理变量。

指标有效性的验证如下。时间趋势上，DU\_kw和DU\_llm在2011至2024年间均呈上升趋势，在政策时点加速明显。行业分布上，信息技术产业和科学研究与技术服务业的指标均值远高于制造业和采矿业，符合产业间数据依赖程度的差异。两套测度的Pearson相关系数为0.678，一致性较强但并不重叠：DU\_kw侧重利用的广度，DU\_llm侧重利用的深度。本文进一步通过篇幅控制、去泛化词口径和GenericNarr联合回归检验构念边界，并增加DUchain\_count、DUclosedloop、DUcore和DUkw\_mda等增强口径，以检验更接近业务嵌入和链条完整性的表述是否同样能够预测股价延迟。需要说明的是，受限于现有归档条件，本文暂未能重建基于mean score和share(score\textgreater=2)的全样本替代聚合；但本文已对120条分层上下文样本完成独立人工复核，并报告人工评分与独立Codex样本级评分的一致性统计，以补强语义评分规则的可复核性。

\subsection{（三）变量定义}\label{ux4e09ux53d8ux91cfux5b9aux4e49}

因变量为股价延迟PriceDelay，采用Hou和Moskowitz（2005）的方法构造。对企业\(i\)在年份\(t\)的日收益率序列，要求至少120个交易日，分别估计受限模型和无限制模型：

受限模型： \[r_{i,d} = \alpha + \beta_0 \cdot r_{m,d} + \varepsilon_{i,d} \qquad (2)\]

无限制模型： \[r_{i,d} = \alpha + \beta_0 \cdot r_{m,d} + \sum_{k=1}^{5} \beta_k \cdot r_{m,d-k} + \varepsilon_{i,d} \qquad (3)\]

股价延迟定义为： \[\text{PriceDelay}_{i,t} = 1 - \frac{R^2_{\text{restricted}}}{R^2_{\text{unrestricted}}} \qquad (4)\]

取值范围为{[}0,1{]}，值越大表示股价对市场信息的反应越滞后，定价效率越低。若H1成立，DU应与PriceDelay呈负向关联。

作为辅助指标，本文同时采用股价同步性SYNCH = ln(R²/(1-R²))，其中R²来自个股收益率对市场收益率和行业收益率的回归。PriceDelay衡量价格对市场信息的吸收速度，SYNCH反映市场与行业共同因素对个股收益率的解释比例。若数据要素利用改善了信息环境，SYNCH上升可能反映系统性信息定价更快或特质噪声下降，但其经济含义不同于PriceDelay，因此本文仅将其作为辅助结果报告，而不将其作为H1的镜像验证。

自变量为数据要素利用。主指标DU\_kw和语义增强指标DU\_llm的定义和构造方法已在上一节详述；DUevent、DUchain\_count、DUclosedloop、DUcore和DUkw\_mda则作为支持性增强测度，在后文用于检验``更完整、更嵌入''的数据利用表述是否具有与基准指标一致的定价含量。

控制变量共11个，包括企业规模、财务状况和治理特征：

\textbf{表2 控制变量定义}

{\def\LTcaptype{} % do not increment counter
\begin{longtable}[]{@{}ll@{}}
\toprule\noalign{}
变量 & 定义 \\
\midrule\noalign{}
\endhead
\bottomrule\noalign{}
\endlastfoot
Size & 总市值的自然对数 \\
Lev & 总负债/总资产 \\
ROA & 净利润/总资产 \\
TobinQ & (股权市值+负债账面值)/总资产 \\
Age & ln(当年-上市年份+1) \\
Growth & 营收增长率 \\
IndepRatio & 独立董事人数/董事会总人数 \\
Dual & 董事长兼任CEO取1，否则为0 \\
Top1Share & 第一大股东持股比例 \\
SOE & 国有企业取1，否则为0 \\
CFO & 经营活动现金流/总资产 \\
\end{longtable}
}

控制变量参照股价延迟研究的通行做法选取。由于机制分析关注的是数据要素利用如何影响分析师预测分歧、现金流波动和供应链集中度，这些变量本身不纳入基准控制，以避免``坏控制''问题（Angrist and Pischke, 2009）。

标准误在行业×年份层面聚类，以校正同行业企业在同一年份受共同技术扩散冲击导致的截面相关。稳健性检验中报告了企业和年份双向聚类的结果，结论不变。

\subsection{（四）模型设定}\label{ux56dbux6a21ux578bux8bbeux5b9a}

为检验数据要素利用对资本市场定价效率的影响，构建如下基准回归模型：

\[PriceDelay_{i,t} = \alpha_0 + \alpha_1 DU_{i,t} + \alpha_2 Control_{i,t} + \mu_i + \lambda_t + \varepsilon_{i,t} \qquad (5)\]

其中，PriceDelay为企业\(i\)在第\(t\)年的股价延迟；DU为数据要素利用指标，分别以DU\_kw和DU\_llm进入回归；Control为11个控制变量向量；\(\mu_i\)和\(\lambda_t\)分别为企业固定效应和年份固定效应；\(\varepsilon_{i,t}\)为随机扰动项。标准误在行业×年份层面聚类。若H1成立，\(\alpha_1\)应显著为负。

内生性方面，本文采用两个工具变量：IV1为同行业跨省企业DU的留一均值，利用行业层面技术扩散的外生变异；IV2为省级大数据发展总指数（2016年截面）与年份趋势的交互项，利用地理维度基础设施供给的外生变异。排除限制的论证和诊断检验在第四节报告。Oster界检验、Heckman两阶段法和安慰剂检验作为辅助识别策略。

传导机制的检验模型为：

\[Channel_{i,t} = \beta_0 + \beta_1 DU_{i,t} + \beta_2 Control_{i,t} + \mu_i + \lambda_t + \varepsilon_{i,t} \qquad (6)\]

其中Channel为渠道变量，定义如下：

\textbf{表3 渠道变量定义}

{\def\LTcaptype{} % do not increment counter
\begin{longtable}[]{@{}ll@{}}
\toprule\noalign{}
渠道变量 & 定义 \\
\midrule\noalign{}
\endhead
\bottomrule\noalign{}
\endlastfoot
ForecastDisp 分析师预测分歧 & 同一企业年度内至少3条盈利预测的标准差 \\
CashFlowVol 现金流波动率 & 经营活动现金流/总资产的三年滚动标准差 \\
SCConc 综合供应链集中度 & 前五大客户销售占比与前五大供应商采购占比的均值 \\
\end{longtable}
}

遵循江艇（2022）的建议，采用路径系数的直接呈现来识别传导渠道的存在性，同时使用DU\_kw和DU\_llm两套测度进行交叉验证。双测度下方向一致且均显著的渠道视为更稳健的路径证据；仅在一套测度下显著的渠道一并报告，但对其解释保持审慎。渠道变量均经1\%/99\%缩尾处理。

异质性分析方面，主文在统一滞后规格下优先报告三条仍具有较清晰解释力的边界条件：政策阶段差异、高分析师关注与非主板企业。具体地，政策阶段以2022年为切点；高分析师关注按分析师覆盖的全样本中位数分组；MainBoard取1表示主板公司、0表示非主板公司。数字经济核心产业、替代切点以及其他补充分组结果仍放在附录中，作为补充证据报告。考虑到统一滞后规格后部分产业差异不再稳健，本文对异质性结论主要采用交互项检验；政策阶段结果另辅以费舍尔组合检验并使用更审慎的解释口径。
